% TOP_PROBLEM: {"id": 3335, "title": "Tarski's exponential function problem", "category_id": 18, "category": {"id": 18, "name": "logic", "display_name": "Logic", "description": "Problems in mathematical logic, model theory, proof theory, and finite model theory.", "slug": "logic", "order_index": 18, "created_at": "2026-07-31T15:26:25.671Z"}, "status": "open", "problem_number": "OPG-1790", "set_id": 12}
% FIRST_POSTED: 2026-10-01
% ENTRY_KIND: statement_only
% UnsolvedMath ID 3335; source catalogue code OPG-1790
% !TeX program = lualatex
% Definition and sources only; this file is not an LLM solution attempt.
\documentclass[11pt,a4paper]{article}
\usepackage[margin=25mm]{geometry}
\usepackage{fontspec}
\setmainfont{lmroman10-regular.otf}[BoldFont=lmroman10-bold.otf,
 ItalicFont=lmroman10-italic.otf,BoldItalicFont=lmroman10-bolditalic.otf]
\usepackage{amsmath,amssymb,mathtools}
\usepackage{unicode-math}
\setmathfont{latinmodern-math.otf}
\AtBeginDocument{\let\setminus\smallsetminus}
\usepackage{xurl,hyperref}
\hypersetup{colorlinks=true,urlcolor=blue}
\setcounter{secnumdepth}{0}
\setlength{\parindent}{0pt}
\setlength{\parskip}{5pt}
\setlength{\emergencystretch}{3em}
\begin{document}
\section*{Tarski's exponential function problem}\label{3335}
{\small UnsolvedMath ID 3335 · OPG-1790 · Logic · OpenGarden}
\subsection{Definitions and mathematical statement}
Consider the ordered real exponential field
\[
             \mathbb R_{\exp}=(\mathbb R;0,1,+,\cdot,<,\exp),
             \qquad \exp(x)=e^x.
\]
Its first-order language has the displayed constant, function, and relation symbols, together with equality, variables, Boolean connectives, and quantifiers over individual real numbers. Terms are built by finitely many applications of addition, multiplication, and the total exponential function. Quantification over subsets, functions, or integers as a separate domain is not part of this language.

The complete first-order theory $\operatorname{Th}(\mathbb R_{\exp})$ is the set of all sentences in this language that are true under the usual real interpretation. The decision problem asks whether there is an algorithm which, given any finite sentence, always halts and correctly decides whether it belongs to this theory.

The algorithm must handle arbitrary finite alternations of existential and universal quantifiers and arbitrary nesting of exponentials. There are no arbitrary real-number parameters in the input: coefficients are specified using the named constants and finite expressions, with rational constants available through their usual definitions. Allowing unrestricted named real constants without effective encodings would be a different problem.

Tarski's exponential-function question asks for unconditional decidability or a proof of undecidability of this exact theory. Deciding only quantifier-free instances, only bounded numerical approximations, or a restricted exponential function on a compact interval would not settle it. A procedure justified under an additional number-theoretic conjecture would likewise be a conditional answer rather than the requested unconditional one.
\subsection{Short English statement}
Is there an always-terminating algorithm deciding every first-order sentence about the real field with its total exponential function?
\subsection{Sources}
\begin{itemize}
\item[S1] ULAM.AI and the UnsolvedMath Contributors, supplied problems.json, record 3335 (OPG-1790), OpenGarden. Catalogue identity and source transcription; CC BY 4.0.
\url{https://huggingface.co/datasets/ulamai/UnsolvedMath}.
\item[S2] Open Problem Garden, Tarski's exponential function problem. Original problem and discussion.
\url{http://www.openproblemgarden.org/op/tarskis_exponential_function_problem}.
\item[S3] A. Berarducci and F. Gallinaro, On the elementary theory of the real exponential field (2026).
\url{https://arxiv.org/abs/2603.08365}.
\end{itemize}
\end{document}
